\documentclass[10pt,a4paper]{amsart}
\usepackage{amsmath,amssymb,amsthm,graphicx}
% Retained preamble for the Stacks Project chapter morphisms.tex (Morphisms of Schemes).
% The chapter file pinned-original-source.tex carries only \input{preamble} on its first line, so no
% preamble text lives inside the chapter itself.  The Stacks Project uses one single shared preamble for
% every chapter; the snapshot stages it verbatim in sources/stacks-chow/preamble.tex.  The block below is
% copied mechanically, line for line, from that file's lines 14-18 (the Xy-pic setup) and lines 154-217
% (the native theorem declarations together with the mathematical macros).  Nothing was added, removed or
% redefined.  The wrapper class and the amsmath/amssymb/amsthm/graphicx base packages are supplied by the
% assembler, so the class/\externaldocument part of the shared preamble is not repeated here.
% For commutative diagrams we use Xy-pic
\usepackage[all]{xy}

% We use 2cell for 2-commutative diagrams.
\xyoption{2cell}
% Theorem environments.
%
\theoremstyle{plain}
\newtheorem{theorem}[subsection]{Theorem}
\newtheorem{proposition}[subsection]{Proposition}
\newtheorem{lemma}[subsection]{Lemma}

\theoremstyle{definition}
\newtheorem{definition}[subsection]{Definition}
\newtheorem{example}[subsection]{Example}
\newtheorem{exercise}[subsection]{Exercise}
\newtheorem{situation}[subsection]{Situation}

\theoremstyle{remark}
\newtheorem{remark}[subsection]{Remark}
\newtheorem{remarks}[subsection]{Remarks}

\numberwithin{equation}{subsection}

% Macros
%
\def\lim{\mathop{\mathrm{lim}}\nolimits}
\def\colim{\mathop{\mathrm{colim}}\nolimits}
\def\Spec{\mathop{\mathrm{Spec}}}
\def\Hom{\mathop{\mathrm{Hom}}\nolimits}
\def\Ext{\mathop{\mathrm{Ext}}\nolimits}
\def\SheafHom{\mathop{\mathcal{H}\!\mathit{om}}\nolimits}
\def\SheafExt{\mathop{\mathcal{E}\!\mathit{xt}}\nolimits}
\def\Sch{\mathit{Sch}}
\def\Mor{\mathop{\mathrm{Mor}}\nolimits}
\def\Ob{\mathop{\mathrm{Ob}}\nolimits}
\def\Sh{\mathop{\mathit{Sh}}\nolimits}
\def\NL{\mathop{N\!L}\nolimits}
\def\CH{\mathop{\mathrm{CH}}\nolimits}
\def\proetale{{pro\text{-}\acute{e}tale}}
\def\etale{{\acute{e}tale}}
\def\QCoh{\mathit{QCoh}}
\def\Ker{\mathop{\mathrm{Ker}}}
\def\Im{\mathop{\mathrm{Im}}}
\def\Coker{\mathop{\mathrm{Coker}}}
\def\Coim{\mathop{\mathrm{Coim}}}

% Boxtimes
%
\DeclareMathSymbol{\boxtimes}{\mathbin}{AMSa}{"02}

%
% Macros for moduli stacks/spaces
%
\def\QCohstack{\mathcal{QC}\!\mathit{oh}}
\def\Cohstack{\mathcal{C}\!\mathit{oh}}
\def\Spacesstack{\mathcal{S}\!\mathit{paces}}
\def\Quotfunctor{\mathrm{Quot}}
\def\Hilbfunctor{\mathrm{Hilb}}
\def\Curvesstack{\mathcal{C}\!\mathit{urves}}
\def\Polarizedstack{\mathcal{P}\!\mathit{olarized}}
\def\Complexesstack{\mathcal{C}\!\mathit{omplexes}}
% \Pic is the operator that assigns to X its picard group, usage \Pic(X)
% \Picardstack_{X/B} denotes the Picard stack of X over B
% \Picardfunctor_{X/B} denotes the Picard functor of X over B
\def\Pic{\mathop{\mathrm{Pic}}\nolimits}
\def\Picardstack{\mathcal{P}\!\mathit{ic}}
\def\Picardfunctor{\mathrm{Pic}}
\def\Deformationcategory{\mathcal{D}\!\mathit{ef}}

% --- appended 2026-10-12 (Chief): the project own definitions of the reference/slogan/history
% display environments, copied mechanically from the shared Stacks preamble lines 8-12
% (\usepackage{verbatim} plus the three \newenvironment lines, which the chapter preamble keeps
% verbatim from the project preamble). Nothing added, removed or redefined: these environments
% carry no mathematical content and the project itself typesets them as comments. Without these
% three lines a counted statement whose own text uses them fails with Environment slogan undefined.
\usepackage{verbatim}
\newenvironment{reference}{\comment}{\endcomment}
\newenvironment{slogan}{\comment}{\endcomment}
\newenvironment{history}{\comment}{\endcomment}
\usepackage{cleveref}
\crefname{originallemma}{original lemma}{original lemmas}
\Crefname{originallemma}{Original lemma}{Original lemmas}
\makeatletter
\newlabel{algebra-lemma-dominate}{{00IA}{source}{}{sourcepointer-1}{}}
\newlabel{algebra-lemma-dominate@cref}{{[originallemma][0][]00IA}{[1][0][]source}{}{}{}}
\makeatother
\crefname{originallemma}{original lemma}{original lemmas}
\Crefname{originallemma}{Original lemma}{Original lemmas}
\makeatletter
\newlabel{algebra-lemma-stack-valuation-rings}{{088Z}{source}{}{sourcepointer-2}{}}
\newlabel{algebra-lemma-stack-valuation-rings@cref}{{[originallemma][0][]088Z}{[1][0][]source}{}{}{}}
\makeatother
\crefname{originallemma}{original lemma}{original lemmas}
\Crefname{originallemma}{Original lemma}{Original lemmas}
\makeatletter
\newlabel{lemma-flat-base-change-scheme-theoretic-image}{{081I}{source}{}{sourcepointer-3}{}}
\newlabel{lemma-flat-base-change-scheme-theoretic-image@cref}{{[originallemma][0][]081I}{[1][0][]source}{}{}{}}
\makeatother
\crefname{originallemma}{original lemma}{original lemmas}
\Crefname{originallemma}{Original lemma}{Original lemmas}
\makeatletter
\newlabel{lemma-reach-points-scheme-theoretic-image}{{02JQ}{source}{}{sourcepointer-4}{}}
\newlabel{lemma-reach-points-scheme-theoretic-image@cref}{{[originallemma][0][]02JQ}{[1][0][]source}{}{}{}}
\makeatother
\crefname{originallemma}{original lemma}{original lemmas}
\Crefname{originallemma}{Original lemma}{Original lemmas}
\makeatletter
\newlabel{lemma-scheme-theoretic-image-reduced}{{056B}{source}{}{sourcepointer-5}{}}
\newlabel{lemma-scheme-theoretic-image-reduced@cref}{{[originallemma][0][]056B}{[1][0][]source}{}{}{}}
\makeatother
\crefname{originallemma}{original lemma}{original lemmas}
\Crefname{originallemma}{Original lemma}{Original lemmas}
\makeatletter
\newlabel{schemes-lemma-map-into-reduction}{{0356}{source}{}{sourcepointer-6}{}}
\newlabel{schemes-lemma-map-into-reduction@cref}{{[originallemma][0][]0356}{[1][0][]source}{}{}{}}
\makeatother
\crefname{originalproposition}{original proposition}{original propositions}
\Crefname{originalproposition}{Original proposition}{Original propositions}
\makeatletter
\newlabel{schemes-proposition-characterize-universally-closed}{{01KF}{source}{}{sourcepointer-7}{}}
\newlabel{schemes-proposition-characterize-universally-closed@cref}{{[originalproposition][0][]01KF}{[1][0][]source}{}{}{}}
\makeatother
\title{A refined valuative criterion with a dense subscheme implies universal closedness, and separatedness when the morphism is quasi-separated}
\author{The Stacks Project Authors}
\date{stacks-project:morphisms.tex:local-snapshot-20261012; GNU FDL 1.2 or later}
\begin{document}
\maketitle
\noindent\textit{Editorial note.} All delivered mathematical text is the author's own, copied from the pinned Stacks Project chapter file pinned-original-source.tex; no statement or proof step was rewritten, reordered or removed.\par\medskip
\noindent\textit{Editorial note.} Following the multi-claim convention for this source, source\_url is the source-level URL https://stacks.math.columbia.edu/ and the source label lemma-refined-valuative-criterion-universally-closed is carried in claim\_key and here rather than in a per-claim tag URL.\par\medskip
\noindent\textit{Editorial note.} The preamble is the chapter's shared Stacks Project preamble, retained as sources/stacks-morphisms/retained-preamble.tex: a line-for-line excerpt of the shared Stacks preamble staged with this snapshot. No macro was invented or redefined, and this card adds no preamble text of its own.\par\medskip
\noindent\textit{Editorial note.} Statement span 10219-10234 and proof span 10236-10327 come from the halfA pool card operations/pool-cards/stacks-morphisms-halfA.json and were re-read in the pinned chapter file before this card was written.\par\medskip
\noindent\textit{Editorial note.} The statement and proof contain four xymatrix solid diagrams with a dotted arrow; they are copied verbatim and the retained chapter preamble loads Xy-pic.\par\medskip
\noindent\textit{Editorial note.} The proof uses the dominance relation on valuation rings inside a field exactly as the author writes it, including the conclusion that the given valuation ring equals the intersection of the field with the ambient residue field.\par\medskip
\noindent\textit{Editorial note.} 7 cross-references occur inside the counted spans; they are listed in external\_references as navigation pointers only.\par\medskip
\noindent\textit{Allowed standard prerequisites (editorial).}\begin{itemize}
\item The characterization of universally closed morphisms by the existence part of the valuative criterion, cited from the schemes chapter. Cited at proof lines 10237-10239. Reduces universal closedness of the given morphism to the existence of the dotted arrow in a valuation ring diagram.
\item Reductions of schemes and the behaviour of morphisms and dense images under passage to the reduction, together with the description of the scheme theoretic image in the reduced case, cited from the schemes chapter and the same chapter. Cited at proof lines 10244-10249. Lets the proof assume all schemes reduced and rewrite density of the image as the scheme theoretic image being the whole space.
\item Flatness of a polynomial algebra over the coordinate ring of an affine open and invariance of the scheme theoretic image under flat base change, cited from the same chapter. Cited at proof lines 10256-10259. Shows that the scheme theoretic image of the base changed dense open is the whole of the polynomial algebra's spectrum.
\item Existence of a valuation ring diagram reaching a prescribed point of a scheme theoretic image, cited from the same chapter. Cited at proof lines 10260-10262. Produces the valuation ring A' and the field K' whose closed point maps to the point of interest.
\item Existence of a valuation ring inside a residue field dominating the image of a given valuation ring, and the description of a valuation ring sandwiched between a valuation ring and a residue field, both cited from the algebra chapter. Cited at proof lines 10271-10278. Builds the valuation ring A'' and the auxiliary valuation ring C that put the situation into the shape of the hypothesis.
\end{itemize}\medskip
\ifdefined\hypertarget\hypertarget{sourcepointer-1}{}\fi
\noindent\textit{Source reference.} Original lemma 00IA: focused-seeds/AGP/sources/stacks-snapshot-20261011/tags.txt, source lines 672--672. Same-chapter reference used in the retained author text at source line(s) 941, 10279. It is supplied only as a navigation pointer so that the retained text compiles; the counted item delivers the author's own statement and the author's complete proof, and no part of the referenced item is copied in or substituted for the counted argument.\par\medskip
\ifdefined\hypertarget\hypertarget{sourcepointer-2}{}\fi
\noindent\textit{Source reference.} Original lemma 088Z: focused-seeds/AGP/sources/stacks-snapshot-20261011/tags.txt, source lines 10479--10479. Same-chapter reference used in the retained author text at source line(s) 10284. It is supplied only as a navigation pointer so that the retained text compiles; the counted item delivers the author's own statement and the author's complete proof, and no part of the referenced item is copied in or substituted for the counted argument.\par\medskip
\ifdefined\hypertarget\hypertarget{sourcepointer-3}{}\fi
\noindent\textit{Source reference.} Original lemma 081I: focused-seeds/AGP/sources/stacks-snapshot-20261011/tags.txt, source lines 10218--10218. Same-chapter reference used in the retained author text at source line(s) 10264. It is supplied only as a navigation pointer so that the retained text compiles; the counted item delivers the author's own statement and the author's complete proof, and no part of the referenced item is copied in or substituted for the counted argument.\par\medskip
\ifdefined\hypertarget\hypertarget{sourcepointer-4}{}\fi
\noindent\textit{Source reference.} Original lemma 02JQ: focused-seeds/AGP/sources/stacks-snapshot-20261011/tags.txt, source lines 3345--3345. Same-chapter reference used in the retained author text at source line(s) 1367, 4438, 10265, 14326. It is supplied only as a navigation pointer so that the retained text compiles; the counted item delivers the author's own statement and the author's complete proof, and no part of the referenced item is copied in or substituted for the counted argument.\par\medskip
\ifdefined\hypertarget\hypertarget{sourcepointer-5}{}\fi
\noindent\textit{Source reference.} Original lemma 056B: focused-seeds/AGP/sources/stacks-snapshot-20261011/tags.txt, source lines 6704--6704. Same-chapter reference used in the retained author text at source line(s) 1207, 10253. It is supplied only as a navigation pointer so that the retained text compiles; the counted item delivers the author's own statement and the author's complete proof, and no part of the referenced item is copied in or substituted for the counted argument.\par\medskip
\ifdefined\hypertarget\hypertarget{sourcepointer-6}{}\fi
\noindent\textit{Source reference.} Original lemma 0356: focused-seeds/AGP/sources/stacks-snapshot-20261011/tags.txt, source lines 4117--4117. Same-chapter reference used in the retained author text at source line(s) 291, 987, 10250, 13640, 14408, 14563. It is supplied only as a navigation pointer so that the retained text compiles; the counted item delivers the author's own statement and the author's complete proof, and no part of the referenced item is copied in or substituted for the counted argument.\par\medskip
\ifdefined\hypertarget\hypertarget{sourcepointer-7}{}\fi
\noindent\textit{Source reference.} Original proposition 01KF: focused-seeds/AGP/sources/stacks-snapshot-20261011/tags.txt, source lines 2048--2048. Same-chapter reference used in the retained author text at source line(s) 10208, 10239, 10613. It is supplied only as a navigation pointer so that the retained text compiles; the counted item delivers the author's own statement and the author's complete proof, and no part of the referenced item is copied in or substituted for the counted argument.\par\medskip
\begin{lemma}
\label{lemma-refined-valuative-criterion-universally-closed}
Let $f : X \to S$ and $h : U \to X$ be morphisms of schemes.
Assume that $f$ and $h$ are quasi-compact and that $h(U)$ is dense in $X$.
If given any commutative solid diagram
$$
\xymatrix{
\Spec(K) \ar[r] \ar[d] & U \ar[r]^h & X \ar[d]^f \\
\Spec(A) \ar[rr] \ar@{-->}[rru] & & S
}
$$
where $A$ is a valuation ring with field of fractions $K$, there
exists a unique dotted arrow making the diagram commute, then $f$
is universally closed. If moreover $f$ is quasi-separated, then
$f$ is separated.
\end{lemma}

\begin{proof}
To prove $f$ is universally closed we will verify the existence part of the
valuative criterion for $f$ which suffices by
Schemes, Proposition \ref{schemes-proposition-characterize-universally-closed}.
To do this, consider a commutative diagram
$$
\xymatrix{
\Spec(K) \ar[r] \ar[d] & X \ar[d] \\
\Spec(A) \ar[r] & S
}
$$
where $A$ is a valuation ring and $K$ is the fraction field of $A$.
Note that since valuation rings and fields are reduced, we may
replace $U$, $X$, and $S$ by their respective reductions
by Schemes, Lemma \ref{schemes-lemma-map-into-reduction}.
In this case the assumption that $h(U)$ is dense means that
the scheme theoretic image of $h : U \to X$ is $X$, see
Lemma \ref{lemma-scheme-theoretic-image-reduced}.
We may also replace $S$ by an affine open through which
the morphism $\Spec(A) \to S$ factors. Thus we may
assume that $S = \Spec(R)$.

\medskip\noindent
Let $\Spec(B) \subset X$ be an affine open through which
the morphism $\Spec(K) \to X$ factors. Choose a polynomial
algebra $P$ over $B$ and a $B$-algebra surjection $P \to K$.
Then $\Spec(P) \to X$ is flat. Hence the scheme theoretic image
of the morphism $U \times_X \Spec(P) \to \Spec(P)$ is $\Spec(P)$ by
Lemma \ref{lemma-flat-base-change-scheme-theoretic-image}.
By Lemma \ref{lemma-reach-points-scheme-theoretic-image}
we can find a commutative diagram
$$
\xymatrix{
\Spec(K') \ar[r] \ar[d] & U \times_X \Spec(P) \ar[d] \\
\Spec(A') \ar[r] & \Spec(P)
}
$$
where $A'$ is a valuation ring and $K'$ is the fraction field of $A'$
such that the closed point of $\Spec(A')$ maps to $\Spec(K) \subset \Spec(P)$.
In other words, there is a $B$-algebra map
$\varphi : K \to A'/\mathfrak m_{A'}$. Choose a valuation ring
$A'' \subset A'/\mathfrak m_{A'}$ dominating $\varphi(A)$ with
field of fractions $K'' = A'/\mathfrak m_{A'}$
(Algebra, Lemma \ref{algebra-lemma-dominate}). We set
$$
C = \{\lambda \in A' \mid \lambda \bmod \mathfrak m_{A'} \in A''\}.
$$
which is a valuation ring by
Algebra, Lemma \ref{algebra-lemma-stack-valuation-rings}.
As $C$ is an $R$-algebra with fraction field $K'$, we obtain a
commutative diagram
$$
\xymatrix{
\Spec(K') \ar[r] \ar[d] & U \ar[r] & X \ar[d] \\
\Spec(C) \ar[rr] \ar@{-->}[rru] & & S
}
$$
as in the statement of the lemma. Thus a dotted arrow fitting into
the diagram as indicated. By the uniqueness assumption of the lemma
the composition $\Spec(A') \to \Spec(C) \to X$ agrees with the
given morphism $\Spec(A') \to \Spec(P) \to \Spec(B) \subset X$.
Hence the restriction of the morphism to the spectrum of
$C/\mathfrak m_{A'} = A''$ induces the given morphism
$\Spec(K'') = \Spec(A'/\mathfrak m_{A'}) \to \Spec(K) \to X$.
Let $x \in X$ be the image of the closed point of $\Spec(A'') \to X$.
The image of the induced ring map $\mathcal{O}_{X, x} \to A''$
is a local subring which is contained in $K \subset K''$.
Since $A$ is maximal for the relation of domination in $K$
and since $A \subset A''$, we have $A = K \cap A''$. We conclude
that $\mathcal{O}_{X, x} \to A''$ factors through $A \subset A''$.
In this way we obtain our desired arrow $\Spec(A) \to X$.

\medskip\noindent
Finally, assume $f$ is quasi-separated. Then $\Delta : X \to X \times_S X$
is quasi-compact. Given a solid diagram
$$
\xymatrix{
\Spec(K) \ar[r] \ar[d] & U \ar[r]^h & X \ar[d]^\Delta \\
\Spec(A) \ar[rr] \ar@{-->}[rru] & & X \times_S X
}
$$
where $A$ is a valuation ring with field of fractions $K$, there
exists a unique dotted arrow making the diagram commute. Namely,
the lower horizontal arrow is the same thing as a pair of morphisms
$\Spec(A) \to X$ which can serve as the dotted arrow in the diagram
of the lemma. Thus the required uniqueness shows that the lower
horizontal arrow factors through $\Delta$.
Hence we can apply the result we just proved to
$\Delta : X \to X \times_S X$ and $h : U \to X$ and conclude that
$\Delta$ is universally closed. Clearly this means that $f$
is separated.
\end{proof}

\end{document}
